\begin{frame}{Faithful does not imply understandable}
\TradeoffHeader{\ContentIcon}{Content}{Correctness / faithfulness}{\PresentationIcon}{Presentation}{Compactness / composition}{Model-relative validity and human readability are different claims.}
\[
\text{Faithful}\not\Rightarrow\text{Understandable}
\]
A technically faithful explanation with a large number of predicates may be unusable for its audience.
\end{frame}

\begin{frame}{Plausible does not imply faithful}
\TradeoffHeader{\UserIcon}{User}{Coherence}{\ContentIcon}{Content}{Correctness / faithfulness}{Human plausibility is not evidence that the explanation tracks the model.}
\[
\text{Coherent with human knowledge}\not\Rightarrow\text{Correct relative to model}
\]
A domain expert may find an explanation plausible even if the model relies on a spurious feature.
\end{frame}

\begin{frame}{Human-grounded tasks}
\HierarchyHeader{\UserIcon}{User}{Coherence / context / controllability}{Does the explanation improve a controlled human task?}
\begin{itemize}
  \item Predict the model output; detect a model failure.
  \item Identify influential variables; answer comprehension questions.
  \item Measure decision time, confidence, and calibration.
\end{itemize}
Preference alone is weak evidence: ``Which explanation do you like?'' does not test task performance.
\end{frame}
